\clearpage

\lehead[]{\normalfont\sffamily\hspace*{-2.00cm}\textcolor{white}{\colorbox{lightblue}{\parbox[c][0.70cm][b]{1.60cm}{
\makebox[1.60cm][r]{\thechapter}\\ \makebox[1.60cm][r]{ÜBUNG}}}}\hspace{0.17cm}\textcolor{lightblue}{\chaptertitle}}
\rohead[]{\textcolor{lightblue}{\chaptertitle}\normalfont\sffamily\hspace*{0.17cm}\textcolor{white}{\colorbox{lightblue}{\parbox[c][0.70cm][b]{1.60cm}{\thechapter\\
ÜBUNG}}}\hspace{-2.00cm}}
%\chead[]{}
\rehead[]{\textcolor{lightblue}{AvHG, Inf, My}}
\lohead[]{\textcolor{lightblue}{AvHG, Inf, My}}

\section{Übung 1: Einfache Abfragen auf einer Tabelle \tagPflicht}
\noindent \emph{Neue Sprachmittel: \texttt{SELECT}, \texttt{WHERE}, \texttt{DISTINCT}, \texttt{ORDER BY}, Aggregatfunktionen (\texttt{COUNT}, \texttt{AVG}), \texttt{GROUP BY} und \texttt{HAVING}.}

Formuliere für die folgenden Aufgaben jeweils eine passende
\lstinline|SELECT|-Abfrage auf der Tabelle \myUserInput{tier}. Speichere deine
SQL-Anweisungen in einer Textdatei ab.

\begin{compactenum}[a)]
\item Zeige die gesamte Tabelle \myUserInput{tier} an.
\item Zeige nur die Spalten \myUserInput{name} und \myUserInput{tierart} an.
\item Liste die Spalten \myUserInput{name} und \myUserInput{tierart} auf.
Sortiere dabei in aufsteigender alphabetischer Reihenfolge nach dem Namen. An
zweiter Stelle soll nach der Tierart sortiert werden (ebenfalls aufsteigend).
\item Zeige nur die Spalte \myUserInput{geburtstag} an. Sorge dafür, dass jeder
Wert nur einmal angezeigt wird.
\item Zeige die Spalten \myUserInput{name} und \myUserInput{tierart} für alle
Tiere an, die noch leben (\lstinline|lebendig = 'ja'|).
\item Liste die Namen aller Tiere auf, die vor dem Jahr $2004$ geboren wurden
und noch nicht tot sind (\lstinline|todestag IS NULL|). Sortiere die Namen in
absteigender alphabetischer Reihenfolge.
\item Liste alle Spalten der Tiere auf, die weder Hund noch Katze sind.
\item Zeige die Namen und die Geburtstage aller Tiere an, für die kein
Todestag angegeben wurde.
\item Zähle, wie viele Tiere nach dem 01.01.2003 geboren wurden.
\item Liste auf, wie viele Tiere es von jeder Tierart gibt. Das Ergebnis soll
eine Tabelle mit der Anzahl und der jeweiligen Tierart sein.
\item Liste alle Tierarten auf, von denen es zwei oder mehr Tiere gibt.
\item Zähle die Anzahl der unterschiedlichen Tierarten, die in der Tabelle
vorkommen.
\end{compactenum}

